% figures/w05-course-vs-heading.svg — 같은 유도법칙, 두 가지 자동조타.
%
%   bash _tools/tikz2svg.sh figures/src/w05-course-vs-heading.tex figures/w05-course-vs-heading.svg
%
% 사용자 지시 (2026-10-13): 2025년 강의자료 "2. USV 제어기 설계" 슬라이드 60-66 의
% 블록도를 "최대한 비슷하게" 이 강의에 맞춰 다시 그린다. 원본은 Fossen 의 TTK 4190
% 강의노트 그림이므로 **옮겨 붙이지 않고 같은 구조를 직접 그린다** (CLAUDE.md §2).
%
% 이 그림이 말하는 것 하나:
%   LOS 계열의 유도법칙은 **침로(course) 명령**을 낸다. 그런데 이 강의의 배는
%   **선수각(heading) 자동조타**를 쓴다. 그래서 chi_d 를 psi_d 로 바꾸는 자리가
%   필요하고, 그 자리에 들어가는 값이 크랩각 beta 다. beta 를 모르면 그것은
%   요(yaw)에 들어오는 **외란**이고, 5-5 절의 적분이 그 외란을 메운다.
\documentclass[border=6pt]{standalone}
% gnc-style.tex — 이 볼트의 TikZ 그림이 공유하는 스타일.
%
% 저널 그림 규약을 스타일 하나로 굳혀 둔다. 그림마다 선 굵기나 화살촉을
% 다시 정하지 않는다 — 그것이 그림이 제각각으로 보이는 원인이다.
%
% 쓰는 법:  % gnc-style.tex — 이 볼트의 TikZ 그림이 공유하는 스타일.
%
% 저널 그림 규약을 스타일 하나로 굳혀 둔다. 그림마다 선 굵기나 화살촉을
% 다시 정하지 않는다 — 그것이 그림이 제각각으로 보이는 원인이다.
%
% 쓰는 법:  % gnc-style.tex — 이 볼트의 TikZ 그림이 공유하는 스타일.
%
% 저널 그림 규약을 스타일 하나로 굳혀 둔다. 그림마다 선 굵기나 화살촉을
% 다시 정하지 않는다 — 그것이 그림이 제각각으로 보이는 원인이다.
%
% 쓰는 법:  \input{gnc-style}  (figures/src/ 안에서)

\usepackage{tikz}
\usepackage{amsmath}
\usepackage{xcolor}
\usetikzlibrary{arrows.meta, positioning, calc, fit, backgrounds, decorations.pathmorphing}

% ---- 색 --------------------------------------------------------------------
% 색은 강조 하나에만 쓴다. 나머지는 검정.
\definecolor{ink}{HTML}{1A1A1A}
\definecolor{accent}{HTML}{7C3AED}
\definecolor{warn}{HTML}{B91C1C}
\definecolor{soft}{HTML}{666666}

% 기하 그림(좌표계·유도)에서만 쓰는 넷. 블록선도는 위의 흑백 + accent 로 충분하다.
\definecolor{pathc}{HTML}{0D9488}   % 계획 경로 · 항적
\definecolor{errc}{HTML}{C2410C}    % 오차 (cross-track 등)
\definecolor{alng}{HTML}{B45309}    % along-track
\definecolor{flow}{HTML}{0369A1}    % 조류 · 유속

% ---- 선과 화살촉 -----------------------------------------------------------
\tikzset{
  % 신호선. 화살촉은 작고 뾰족한 것 하나뿐이다.
  sig/.style   = {ink, line width=0.5pt, -{Stealth[length=4pt, width=3pt]}},
  wire/.style  = {ink, line width=0.5pt},
  % 강조 경로 — 파선으로 구분한다. 흑백 인쇄에서도 살아야 하기 때문이다.
  asig/.style  = {accent, line width=0.5pt, densely dashed,
                  -{Stealth[length=4pt, width=3pt]}},
  awire/.style = {accent, line width=0.5pt, densely dashed},
  % 블록. 흰 바탕, 직각, 검은 테두리.
  blk/.style   = {draw=ink, line width=0.55pt, fill=white,
                  minimum width=17mm, minimum height=8mm, inner sep=2pt, font=\small},
  % 합산점
  sum/.style   = {draw=ink, line width=0.55pt, fill=white, circle,
                  inner sep=0pt, minimum size=4.6mm, font=\scriptsize},
  asum/.style  = {draw=accent, line width=0.55pt, fill=white, circle,
                  inner sep=0pt, minimum size=4.6mm, font=\scriptsize},
  % 분기점
  dot/.style   = {ink, fill=ink, circle, inner sep=0pt, minimum size=1.5mm},
  % 신호 이름 — 선 위에 작게
  sname/.style = {font=\small, inner sep=1.5pt},
  % 그림 안의 짧은 설명. 그림 안의 글은 최소로 둔다
  note/.style  = {font=\scriptsize, soft, inner sep=1.5pt},
  anote/.style = {font=\scriptsize, accent, inner sep=1.5pt},
  % 축
  axis/.style  = {ink, line width=0.5pt},
  tick/.style  = {font=\scriptsize, soft},
}

% 캡션 한 줄 — 그림 아래 가는 선과 함께
\newcommand{\gnccaption}[2]{%
  \node[anchor=north west, text width=#1, align=left, font=\scriptsize, ink]
        at (current bounding box.south west) {\vspace{2pt}#2};}
  (figures/src/ 안에서)

\usepackage{tikz}
\usepackage{amsmath}
\usepackage{xcolor}
\usetikzlibrary{arrows.meta, positioning, calc, fit, backgrounds, decorations.pathmorphing}

% ---- 색 --------------------------------------------------------------------
% 색은 강조 하나에만 쓴다. 나머지는 검정.
\definecolor{ink}{HTML}{1A1A1A}
\definecolor{accent}{HTML}{7C3AED}
\definecolor{warn}{HTML}{B91C1C}
\definecolor{soft}{HTML}{666666}

% 기하 그림(좌표계·유도)에서만 쓰는 넷. 블록선도는 위의 흑백 + accent 로 충분하다.
\definecolor{pathc}{HTML}{0D9488}   % 계획 경로 · 항적
\definecolor{errc}{HTML}{C2410C}    % 오차 (cross-track 등)
\definecolor{alng}{HTML}{B45309}    % along-track
\definecolor{flow}{HTML}{0369A1}    % 조류 · 유속

% ---- 선과 화살촉 -----------------------------------------------------------
\tikzset{
  % 신호선. 화살촉은 작고 뾰족한 것 하나뿐이다.
  sig/.style   = {ink, line width=0.5pt, -{Stealth[length=4pt, width=3pt]}},
  wire/.style  = {ink, line width=0.5pt},
  % 강조 경로 — 파선으로 구분한다. 흑백 인쇄에서도 살아야 하기 때문이다.
  asig/.style  = {accent, line width=0.5pt, densely dashed,
                  -{Stealth[length=4pt, width=3pt]}},
  awire/.style = {accent, line width=0.5pt, densely dashed},
  % 블록. 흰 바탕, 직각, 검은 테두리.
  blk/.style   = {draw=ink, line width=0.55pt, fill=white,
                  minimum width=17mm, minimum height=8mm, inner sep=2pt, font=\small},
  % 합산점
  sum/.style   = {draw=ink, line width=0.55pt, fill=white, circle,
                  inner sep=0pt, minimum size=4.6mm, font=\scriptsize},
  asum/.style  = {draw=accent, line width=0.55pt, fill=white, circle,
                  inner sep=0pt, minimum size=4.6mm, font=\scriptsize},
  % 분기점
  dot/.style   = {ink, fill=ink, circle, inner sep=0pt, minimum size=1.5mm},
  % 신호 이름 — 선 위에 작게
  sname/.style = {font=\small, inner sep=1.5pt},
  % 그림 안의 짧은 설명. 그림 안의 글은 최소로 둔다
  note/.style  = {font=\scriptsize, soft, inner sep=1.5pt},
  anote/.style = {font=\scriptsize, accent, inner sep=1.5pt},
  % 축
  axis/.style  = {ink, line width=0.5pt},
  tick/.style  = {font=\scriptsize, soft},
}

% 캡션 한 줄 — 그림 아래 가는 선과 함께
\newcommand{\gnccaption}[2]{%
  \node[anchor=north west, text width=#1, align=left, font=\scriptsize, ink]
        at (current bounding box.south west) {\vspace{2pt}#2};}
  (figures/src/ 안에서)

\usepackage{tikz}
\usepackage{amsmath}
\usepackage{xcolor}
\usetikzlibrary{arrows.meta, positioning, calc, fit, backgrounds, decorations.pathmorphing}

% ---- 색 --------------------------------------------------------------------
% 색은 강조 하나에만 쓴다. 나머지는 검정.
\definecolor{ink}{HTML}{1A1A1A}
\definecolor{accent}{HTML}{7C3AED}
\definecolor{warn}{HTML}{B91C1C}
\definecolor{soft}{HTML}{666666}

% 기하 그림(좌표계·유도)에서만 쓰는 넷. 블록선도는 위의 흑백 + accent 로 충분하다.
\definecolor{pathc}{HTML}{0D9488}   % 계획 경로 · 항적
\definecolor{errc}{HTML}{C2410C}    % 오차 (cross-track 등)
\definecolor{alng}{HTML}{B45309}    % along-track
\definecolor{flow}{HTML}{0369A1}    % 조류 · 유속

% ---- 선과 화살촉 -----------------------------------------------------------
\tikzset{
  % 신호선. 화살촉은 작고 뾰족한 것 하나뿐이다.
  sig/.style   = {ink, line width=0.5pt, -{Stealth[length=4pt, width=3pt]}},
  wire/.style  = {ink, line width=0.5pt},
  % 강조 경로 — 파선으로 구분한다. 흑백 인쇄에서도 살아야 하기 때문이다.
  asig/.style  = {accent, line width=0.5pt, densely dashed,
                  -{Stealth[length=4pt, width=3pt]}},
  awire/.style = {accent, line width=0.5pt, densely dashed},
  % 블록. 흰 바탕, 직각, 검은 테두리.
  blk/.style   = {draw=ink, line width=0.55pt, fill=white,
                  minimum width=17mm, minimum height=8mm, inner sep=2pt, font=\small},
  % 합산점
  sum/.style   = {draw=ink, line width=0.55pt, fill=white, circle,
                  inner sep=0pt, minimum size=4.6mm, font=\scriptsize},
  asum/.style  = {draw=accent, line width=0.55pt, fill=white, circle,
                  inner sep=0pt, minimum size=4.6mm, font=\scriptsize},
  % 분기점
  dot/.style   = {ink, fill=ink, circle, inner sep=0pt, minimum size=1.5mm},
  % 신호 이름 — 선 위에 작게
  sname/.style = {font=\small, inner sep=1.5pt},
  % 그림 안의 짧은 설명. 그림 안의 글은 최소로 둔다
  note/.style  = {font=\scriptsize, soft, inner sep=1.5pt},
  anote/.style = {font=\scriptsize, accent, inner sep=1.5pt},
  % 축
  axis/.style  = {ink, line width=0.5pt},
  tick/.style  = {font=\scriptsize, soft},
}

% 캡션 한 줄 — 그림 아래 가는 선과 함께
\newcommand{\gnccaption}[2]{%
  \node[anchor=north west, text width=#1, align=left, font=\scriptsize, ink]
        at (current bounding box.south west) {\vspace{2pt}#2};}


\definecolor{knownc}{HTML}{1D4ED8}
\definecolor{ctrlc}{HTML}{B45309}

\tikzset{
  blk/.style   = {draw=ink, line width=0.8pt, fill=white, align=center,
                  inner sep=4pt, font=\scriptsize, minimum height=9mm},
  guid/.style  = {blk, draw=accent, line width=1.0pt, fill=accent!5},
  auto/.style  = {blk, draw=ctrlc, line width=1.0pt, fill=alng!6},
  craft/.style = {blk, draw=pathc, line width=1.0pt, fill=pathc!10},
  env/.style   = {draw=soft, line width=0.6pt, dashed, rounded corners=2pt, fill=black!2},
  sig/.style   = {ink, line width=0.8pt, -{Stealth[length=4.4pt,width=3.0pt]}},
  cmd/.style   = {accent, line width=1.1pt, -{Stealth[length=4.8pt,width=3.4pt]}},
  fb/.style    = {soft, line width=0.7pt, -{Stealth[length=4.0pt,width=2.8pt]}},
  bus/.style   = {soft, line width=0.7pt},
  dist/.style  = {errc, line width=0.9pt, -{Stealth[length=4.4pt,width=3.0pt]}},
  lbl/.style   = {font=\scriptsize},
  small/.style = {font=\scriptsize, soft},
  capt/.style  = {font=\scriptsize\bfseries, warn},
  band/.style  = {font=\small\bfseries, ink},
  jn/.style    = {draw=ink, line width=0.8pt, circle, inner sep=0pt, minimum size=3.6mm, fill=white},
}

\begin{document}
\begin{tikzpicture}[x=1mm, y=1mm]

% =====================================================================
% 왼쪽 패널 — 침로 제어 / course control
% =====================================================================
\begin{scope}[shift={(0,0)}]
  \node[blk,  text width=15mm] (WP) at (0,0)   {waypoint\\database};
  \node[guid, text width=15mm] (GD) at (26,0)  {LOS\\guidance law};
  \node[auto, text width=15mm] (AP) at (54,0)  {course\\autopilot};
  \node[craft,text width=15mm] (MC) at (80,0)  {marine\\craft};

  \draw[sig] (WP) -- (GD);
  \node[small, anchor=south] at (13,2.4) {$(x_i^n, y_i^n)$};
  \draw[cmd] (GD) -- (AP);
  \node[lbl, accent, anchor=south] at (40,2.4) {$\chi_d$};
  \draw[sig] (AP) -- (MC);
  \draw[sig] (MC.east) -- (92,0);

  %  환경 외란 / the environment
  \draw[dist] (80,16) -- (MC.north);
  \node[small, errc, anchor=south, align=center] at (80,16.6) {waves, wind\\and current};

  %  안쪽 되먹임: 침로는 **대지속도의 방향**이라 GNSS 가 준다
  \draw[bus] (88,0) -- (88,-12) -- (54,-12);
  \fill[soft] (88,0) circle (0.7);
  \draw[fb] (54,-12) -- (AP.south);
  \node[small, anchor=south] at (62,-11.4) {$\chi$ \ (course over ground)};

  %  바깥 되먹임: 위치
  \draw[bus] (88,-12) -- (88,-22) -- (26,-22);
  \draw[fb] (26,-22) -- (GD.south);
  \node[small, anchor=south] at (50,-21.4) {position $(x^n, y^n)$};

  \node[band, anchor=south] at (40,24) {COURSE CONTROL};
  \node[capt, anchor=north, align=center] at (40,-26)
       {the guidance output is used as it stands};
\end{scope}

% =====================================================================
% 오른쪽 패널 — 선수각 제어 (이 강의) / heading control (this course)
% =====================================================================
\begin{scope}[shift={(136,0)}]
  \node[blk,  text width=15mm] (WP2) at (0,0)   {waypoint\\database};
  \node[guid, text width=15mm] (GD2) at (26,0)  {LOS\\guidance law};
  \node[jn]                    (SU)  at (45,0)  {};
  \node[auto, text width=15mm] (AP2) at (62,0)  {heading\\autopilot};
  \node[craft,text width=15mm] (MC2) at (88,0)  {marine\\craft};

  \draw[sig] (WP2) -- (GD2);
  \node[small, anchor=south] at (13,2.4) {$(x_i^n, y_i^n)$};
  \draw[cmd] (GD2) -- (SU);
  \node[lbl, accent, anchor=south] at (38,2.4) {$\chi_d$};
  \draw[cmd] (SU) -- (AP2);
  \node[lbl, accent, anchor=north] at (53.5,-2.2) {$\psi_d$};
  \draw[sig] (AP2) -- (MC2);
  \draw[sig] (MC2.east) -- (100,0);

  %  빼는 것: 크랩각. **부호는 합산점에 둔다** (CLAUDE.md §5)
  \draw[dist] (45,14) -- (SU.north);
  \node[lbl, errc, anchor=west] at (46.8,8.5) {$\beta$};
  \node[small, anchor=east] at (44.0,5.2) {$-$};
  \node[small, errc, anchor=south, align=center] at (45,14.6) {the crab angle,\\\textbf{unknown}};

  \draw[dist] (88,16) -- (MC2.north);
  \node[small, errc, anchor=south, align=center] at (88,16.6) {waves, wind\\and current};

  %  안쪽 되먹임: 선수각은 나침반이 준다 - 속도가 0 이어도 잰다
  \draw[bus] (96,0) -- (96,-12) -- (62,-12);
  \fill[soft] (96,0) circle (0.7);
  \draw[fb] (62,-12) -- (AP2.south);
  \node[small, anchor=south] at (72,-11.4) {$\psi$ \ (compass)};

  \draw[bus] (96,-12) -- (96,-22) -- (26,-22);
  \draw[fb] (26,-22) -- (GD2.south);
  \node[small, anchor=south] at (55,-21.4) {position $(x^n, y^n)$};

  \node[band, anchor=south, ctrlc] at (45,24) {HEADING CONTROL \ --- \ this course};
  \node[capt, anchor=north, align=center] at (45,-26)
       {the output must be converted, and the converter is $\beta$};
\end{scope}

% =====================================================================
% 아래 — 세 줄의 식. 그림의 요점이 이 셋에 있다
% =====================================================================
\begin{scope}[shift={(0,-44)}]
  \node[anchor=north west, font=\small] at (0,0)
       {$\chi = \psi + \beta$};
  \node[small, anchor=north west, text width=56mm] at (0,-7)
       {course is heading plus crab angle: where the vessel \emph{goes} against where it
        \emph{points}};

  \node[anchor=north west, font=\small] at (76,0)
       {$\beta = \arctan\dfrac{v}{u} = \arcsin\dfrac{v}{\sqrt{u^2+v^2}}$};
  \node[small, anchor=north west, text width=58mm] at (76,-12)
       {computable \textbf{only if $v$ is measured}, and singular as the speed goes to zero};

  \node[anchor=north west, font=\small] at (158,0)
       {$\psi_d = \pi_p - \arctan\!\left(K_p y_e^{\,p}\right) - \beta$};
  \node[small, anchor=north west, text width=66mm] at (158,-9)
       {with $\beta$ unknown the last term is missing, and what is missing acts as a
        \textbf{disturbance in yaw} - which is what the integral of §5-5 supplies};
\end{scope}

\end{tikzpicture}
\end{document}
